\documentclass{mem}
\usepackage{natbib}\usepackage{txfonts}\usepackage{balance}
\usepackage{flushend}
\usepackage{graphicx}
\usepackage[breaklinks,pdftex]{hyperref}
\idline{75}{282}
\begin{document}
\def\teff{$T\rm_{eff }$}
\def\kms{$\mathrm {km s}^{-1}$}


\title{
Local Cosmology and Cosmic Distance Scale
}

   \subtitle{Variable stars to investigate our Local Universe}

\author{
G. \,Fiorentino\inst{1} 
\and M. \, Monelli\inst{1,2} \and \bf{not more than three author lines can be included, the remaining will be listed at the end of the paper, before the acknowledgements section.}
%uncomment the following sentence
%\\
%The remaining authors can be found at the end of the paper.
          }

 

\institute{
Istituto Nazionale di Astrofisica --
Osservatorio Astronomico di Roma, via Frascati 33, Monte Porzio Catone, 00078, Italy
\and
 IAC Instituto de Astrof\'isica de Canarias, Calle V\'ia L\'actea sn, 38201 La Laguna, Tenerife, Spain
\email{giuliana.fiorentino@inaf.it}\\
\bf{not more than five institution can be included in this section, the remaining will go at the end of the paper, before the acknowledgements section. If this is the case, please use here the following commented sentence:}
%\\
%The remaining affiliations can be found at the end of the paper.
}

\authorrunning{Fiorentino}


\titlerunning{Variable stars}

\date{Received: XX-XX-XXXX (Day-Month-Year); Accepted: XX-XX-XXXX (Day-Month-Year)}

\abstract{
We discuss the use and importance of pulsating variable stars as population tracers in Local Group galaxies. Among bright variable crossing the classical instability strip, we mostly focus on RR Lyrae stars and Anomalous Cepheids. We discuss their pulsation properties and how it is possible to use them to constrain the evolution and star formation history of the host galaxy. We discuss RR Lyrae stars as tracers of the old population, and how they can be used to trace the accretion history of large galaxies such as the Milky Way and M31, and also the early chemical evolution. Moreover, we show that the frequency of Anomalous Cepheids follows different relations, and therefore trace the intermediate-age star formation. Finally, we discuss the different methods to derive distances and the impact of the Gaia mission. {\bf NOT MORE THAN 1000 characters including spaces are allowed, please double-check before submit the paper.}
\keywords{Stars: Variables: RR Lyrae stars, Anomalous Cepheids -- Dwarf galaxies: Resolved stellar populations }
}
\maketitle{}

\section{Introduction}
The original paper has been modified in order to be used as
demo file and display 
some of the features of mem.cls.

Pulsating variable stars are at the cornerstone of many astrophysical fields: stellar pulsation, stellar evolution, stellar atmospheres, dynamical phenomena, stellar populations, chemical evolution, distance scale (\citealt{bono97b,marconi15,kolenberg10,benko15,feast87,riess21}). As such, they are powerful instruments to understand the evolution of the host galaxy, as they provide direct insights on the age and metallicity distribution of their parent population (\citealt{desomma21a,martinezvazquez16a}).

Pulsating variable stars are therefore essential beacons of the so-called ``Galactic Archaeology'', which uses resolved nearby galaxies in the Local Group (LG) to answer fundamental questions about galaxy formation and evolution \citep{tolstoy09}. According to the Lambda cold dark matter model, dwarf galaxies are relatively simple systems and the first ones to form. They participate in the build-up of larger galaxies and thus they are ideal targets to unravel the physics governing the formation and early evolution of galactic structures and essential probes of the physical ingredients included in current galaxy formation and evolution models (\citealt{dekel86,sawala10,mayer01c}). In this context, LG galaxies offer a fantastic playground for detailed study of the mechanisms driving galaxy evolution, because of the large variety of dwarf galaxies with different morphological and evolutionary properties, and living in different environments. Moreover, they can be resolved into individual stars, for which ages and metallicities can be determined with the help of stellar evolution and stellar atmospheres theory, together with the use of different complementary observational techniques \citep{gallart05}.


%
\begin{figure*}[t!]
\resizebox{\hsize}{!}{\includegraphics[clip=true]{variables_cmd.png}}
\caption{\footnotesize
{\bf Please prefer the use of .png, .jpeg, .eps format to PDF that can give us difficulty in the production phase}
I, V-I Colour magnitude distribution of RRLs, TIICs, ACs and CCs in the LMC. Reddened apparent magnitude are considered. Small points show the location of RRab (black) and RRc (pink) stars. Yellow and green filled triangles are used for fundamental and first overtone ACs, respectively. Similarly, orange and teal filled small circles show fundamental and first overtone CCs. Large open symbols are used for TIICs: BL Her (grey circles), W Vir (red triangles), and RV Tau (blue circles).
}
\label{fig:cmdlmc}
\end{figure*}
%



\section{Observational properties}

The first discoveries of RRL stars outside the Milky Way (MW) dates back to \citet{baade39}, who reported the discovery of 40 variable stars in the Sculptor dwarf spheroidal (dSph) galaxy. However, it was only in the 50s when a significant number (200) of discoveries occurred \citep{thackeray50}. Interestingly, RRL stars in the closer Magellanic Clouds were found later, and by searching in their globular cluster system (\citealt{thackeray51, thackeray53, thackeray58, alexander60}). While these first works  (\citealt{baade61, vanagt67,swope67,vanagt68,hodge78, nemec88}) benefited from the wide-field provided by photographic data, they suffered from strong limitations on the limit magnitude. More modern instrumentation (digital sensor such as charged-coupled devices) allowed to overcome this problem, but with the price of a small area coverage. Variable stars have been searched for and investigated in all bright MW satellites (\citealt{nemec88,siegel00,bersier02,dallora03,baldacci05,kinemuchi08,coppola15,martinezvazquez16a}). In particular, the Magellanic Clouds have been targeted by microlensing project such as MACHO \citep{alcock96a}, the subsequent implementations of the OGLE project (\citealt{udalski92,soszynski19a}) and, moving to the NIR regime, by the VMC project (\citet{cioni11}). Also, since the beginning of the 2000s and thanks to large survey programs such as the SDSS (Sloan Digital Sky Survey) and DES (Dark Energy Survey), the number of faint MW satellites have dramatically increased, extending the range of properties to very low luminosity (-6 $<$ M$_V <$ 0 mag) low metallicity ([Fe/H]$\sim$ -2). Interestingly, a significant fraction of these systems host a few RRL stars, despite the low mass \citep{vivas20b}.


\begin{table*}
\caption{List of the properties of dwarf galaxies.}
\label{tab:tab1}
\begin{center}
\begin{tabular}{lccccccc}
\hline
\\
Galaxy & [Fe/H] & distance & $N_{RR}$ & $N_{Cep}$ & $N_{ACep}$ &  &  \\
\hline
\\
Carina  & -2.1 &  &  & &  &  & \\
Fornax  & & & & &  &  & \\
Sculptor  & & & & &  &  & \\
Draco   & & & & &  &  & \\
\\
\hline
\end{tabular}
\end{center}
\end{table*}


Table \ref{tab:tab1} lists galaxies in and next to the LG, within 2 Mpc, with known variable stars\footnote{We only exclude few isolated dwarf irregular (dIrr) galaxies for which published data are typically sparse, old ground-based data for few bright Cepheids, such as WLM, Pegasus, Sextans A.}. The table presents name, coordinates, the subgroup of each galaxy, and the number of RRab, RRc, Anomalous, Classic, and Type II Cepheids. The last two columns show the photometric system in which the variable stars data have been collected, and the relevant references.
To summarize the observational properties of bright pulsating variable stars in the IS, we show the OGLE sample of LMC variables in both the reddened CMD (see Figure~\ref{fig:cmdlmc}) and in the Wesenheit diagram [$I-1.54(V-I)$, which is by construction reddening-free\footnote{The Wesenheit is pseudo-magnitude defined as the different between a magnitude in a specific passband $X$ minus a colour term multiplied by coefficient representing the ratio between the selective absoprtion in the $X$ band and the colour excess in the assumed colour.}. See the caption for the details about the colour code adopted to display the different classes and pulsation modes. The data are taken from the OGLE III project, in particular RRLs are presented in \citet{soszynski09a}, CCs in \citet{soszynski08b} and finally TIICs and ACs are discussed in \citet{soszynski08c}. In the CMD the considered variable stars partially overlap, making the separation among the different classes difficult. Since also the period ranges partially overlap, the Wesenheit diagram seems to be the best tool to properly classify different types of variables according to their magnitude/colour/period distribution. This is shown in Figure \ref{fig:per_dist}, which shows the period distribution of the four different types of pulsators, for the same stars as in Figure \ref{fig:cmdlmc}.

    \subsection{RR Lyrae stars}\label{sec:rr_obs}

RRLs stars populate the intersection of the IS with the Horizontal Branch (HB). They are therefore low-mass (0.6-0.8M$_{\odot}$), relatively bright ($M_V \sim +0.5$mag) stars burning He in their core. The RRL IS is well defined in colour, therefore in temperature, and limited to $B$--$V$ colour between $\sim$0.2 and $\sim$0.8 mag \citep{walker89}. The majority of RRL stars are known to pulsate either in the fundamental (\textbf{FU or} RRab) or in the first-overtone (\textbf{FO or} RRc) pulsation mode. The former are characterized by longer period ($\sim$0.45 to $\sim$1 d), larger maximum amplitude (up to above 1 mag in the blue optical pass-bands), and asymmetric, saw-tooth-like light curves. RRc present on the other hand shorter periods ($\sim$0.25 to $\sim$0.45 d), amplitude barely reaching 0.7 mag in the visual, smoother and more sinusoidal light curves. 
Moreover, a fraction of RRL stars that ranges from few percents to 10\%, depending on the mean metallicity of the host system, have been found to pulsate in both the fundamental and the first-overtone modes simultaneously (double pulsators or RRd). They typically have periods close to 0.4 d, small amplitude, and noisy light curves resulting from the superposition of the two modes \citep{braga21}.

RRLs have been observed in virtually all the stellar systems they have been searched for: in globular clusters, in the MW  (halo, bulge, thick disk) and in the M31 (halo, disk) field, as well as in dwarf galaxies, with few rare exceptions such as very small mass spheroidal galaxies in the so-called ultra-faint regime \citep{vivas19}. For this reason, and given that they are primary distance indicators, they provide a complementary, population II route to calibrate the distance scale to estimate the Hubble constant, such as that investigated by the Carnegie-Chicago Hubble Program \citep{beaton16}.




\begin{figure}
\resizebox{\hsize}{!}{\includegraphics[clip=true]{per_distribution.png}}
\caption{
\footnotesize
Period distribution of different types of pulsators. The data and the colour-code are the same as in Figure~\ref{fig:cmdlmc}.}
\label{fig:per_dist}
\end{figure}




\section{Conclusions}

The field of pulsating variable stars is now entering a new era, with new scientific and technological challenges.

\textsc{Ongoing and future surveys -} Large data sets from ongoing surveys allow for continuous discovery of new variable stars. Bigger surveys imply more complex data management and analysis, already well beyond the limits of the human lifetime workload capabilities. Considering the variable stars works in LG galaxies revised here, we see that most of them were small-scale, handcrafted analysis, checking candidates by eye one by one. The present and future challenges have shifted to the industrial era of automatization of the main phases of the process: identification (aiming at the highest completeness with the minimum contamination), characterization (period, amplitude, mean magnitude), and classification of variable sources. A huge effort is now required to develop analysis technique that, making use of artificial intelligence and machine learning algorithms in a big data approach, can deal with the enormous amount of incoming data (e.g. \citealt{richards11,elorrieta16,hernitschek16,holl18,hosenie19}). The best example is the Rubin-LSST which, in the near future will produce of the order of 15 Tb of data every night, and will perform real-time analysis in order to deliver alerts minutes after the observations, for prompt follow-up of varying sources. 

\textsc{Tracing the MW components and its surroundings -} Gaia DR3 will publish time series data for millions of stars over the whole sky. Combined with radial velocities and proper motions, this will allow us to identify the remnants of the MW building blocks, characterize their evolution, and associate their GCs and variable stars population so to have a complete view of the MW building blocks (\citealt{massari19,gallart19}). Gaia's variable stars will also allow us to study at least the closest dwarf satellites, tracing their extension, internal structure, and possible interactions with the MW as traced by streams or extra-tidal features. In the next years, Rubin-LSST will allow us to identify RRL stars out to 400 kpc, thus including the whole MW halo and all the southern MW satellites, with 6-bands single-epoch photometry 3 magnitude fainter than, and therefore highly complementary to, Gaia. Moving to the NIR, the high spatial resolution NIR capabilities of ERIS at the VLT will be able, from 2022, to open a new path to the extremely crowded and reddened regions of the galactic bulge, helping to unravel the early history of the MW. On the other hand, wide-field space facilities such as the Nancy Grace Roman Space Telescope will allow to study RRL stars in the NIR well beyond the LG limits.

\textsc{Anchoring the first step of the distance scale - } The improved parallaxes of Gaia will nail down the systematic errors affecting the first step of the distance scale of both Population I and Population II tracers, as discussed above. Nevertheless, Gaia information alone is not enough, as accurate metallicities for the calibrator stars are necessary. At present, a limited set of stars with accurate metallicities is available, but ongoing projects are rapidly increasing the statistics \citep{crestani21b}. In the near future, new generation instruments such as WEAVE, which will start operating in the first half of 2022, and 4MOST, will allow to substantially increase the number of stars with accurate chemical abundance determinations, and this will be crucial to estimate the metallicity dependence of the distance relations of both RRL stars and Cepheid stars.

%UNCOMMENT if you have more than 5 affiliation in the first page of the paper.
%\noindent {\bf{Authors}}\par
%\noindent {\bf{Authors}}\par
%G. Author^{11,13}$,
%P. Author^{7}$\\


%UNCOMMENT if you have more than 5 affiliation in the first page of the paper.
%\noindent {\bf{Affiliations}}\par
%$^{6}$ Istituto Nazionale di Astrofisica Osservatorio Astronomico di Trieste, Via Tiepolo 11, I-34131 Trieste, Italy\par
%$^{7}$ Istituto Nazionale di Astrofisica - Istituto di Astrofisica e Planetologia Spaziali di Roma, Via del Fosso del Cavaliere, 100, 00133 Roma RM\par
%$^{8}$ ASI-Unit\`a di Ricerca Scientifica, Via del Politecnico snc, 00133 Roma, Italy\par
%$^{9}$ Istituto Nazionale di Astrofisica - Osservatorio Astronomico di Brera, Via E. Bianchi 46, 23807 Merate (LC), Italy\par
%$^{10}$ Max-Planck-Institut f\"ur Extraterrestrische Physik, Giessenbachstrasse, 85748 Garching, Germany\par
%$^{11}$ Dipartimento di Fisica, Universit\`a di Roma Tor Vergata, Via della Ricerca Scientifica 1, I-00133 Roma, Italy\\


 
\begin{acknowledgements}
I am grateful to Francesco Primo Bono for comments and suggestions that helped the readability of the paper.
\end{acknowledgements}


\begin{thebibliography}{51}
\bibitem[{{Alcock} {et~al.}(1996){Alcock}, {Allsman}, {Axelrod}, {Bennett},
  {Cook}, {Freeman}, {Griest}, {Marshall}, {Peterson}, {Pratt}, {Quinn},
  {Rodgers}, {Stubbs}, {Sutherland}, \& {Welch}}]{alcock96a}
{Alcock}, C., {Allsman}, R.~A., {Axelrod}, T.~S., {et~al.} 1996, \aj, 111, 1146

\bibitem[{{Alexander}(1960)}]{alexander60}
{Alexander}, J.~B. 1960, \mnras, 121, 97

\bibitem[{{Baade} \& {Hubble}(1939)}]{baade39}
{Baade}, W. \& {Hubble}, E. 1939, \pasp, 51, 40

\bibitem[{{Baade} \& {Swope}(1961)}]{baade61}
{Baade}, W. \& {Swope}, H.~H. 1961, \aj, 66, 300

\bibitem[{{Baldacci} {et~al.}(2005){Baldacci}, {Rizzi}, {Clementini}, \&
  {Held}}]{baldacci05}
{Baldacci}, L., {Rizzi}, L., {Clementini}, G., \& {Held}, E.~V. 2005, \aap,
  431, 1189

\bibitem[{{Beaton} {et~al.}(2016){Beaton}, {Freedman}, {Madore}, {Bono},
  {Carlson}, {Clementini}, {Durbin}, {Garofalo}, {Hatt}, {Jang}, {Kollmeier},
  {Lee}, {Monson}, {Rich}, {Scowcroft}, {Seibert}, {Sturch}, \&
  {Yang}}]{beaton16}
{Beaton}, R.~L., {Freedman}, W.~L., {Madore}, B.~F., {et~al.} 2016, \apj, 832,
  210

\bibitem[{{Benk{\H{o}}} \& {Szab{\'o}}(2015)}]{benko15}
{Benk{\H{o}}}, J.~M. \& {Szab{\'o}}, R. 2015, \apjl, 809, L19

\bibitem[{{Bersier} \& {Wood}(2002)}]{bersier02}
{Bersier}, D. \& {Wood}, P.~R. 2002, \aj, 123, 840

\bibitem[{{Bono} {et~al.}(1997){Bono}, {Caputo}, {Castellani}, \&
  {Marconi}}]{bono97b}
{Bono}, G., {Caputo}, F., {Castellani}, V., \& {Marconi}, M. 1997, \aaps, 121,
  327

\bibitem[{{Braga} {et~al.}(2021){Braga}, {Crestani}, {Fabrizio}, {Bono},
  {Sneden}, {Preston}, {Storm}, {Kamann}, {Latour}, {Lala}, {Lemasle},
  {Prudil}, {Altavilla}, {Chaboyer}, {Dall'Ora}, {Ferraro}, {Gilligan},
  {Fiorentino}, {Iannicola}, {Inno}, {Kwak}, {Marengo}, {Marinoni}, {Marrese},
  {Mart{\'\i}nez-V{\'a}zquez}, {Monelli}, {Mullen}, {Matsunaga}, {Neeley},
  {Stetson}, {Valenti}, \& {Zoccali}}]{braga21}
{Braga}, V.~F., {Crestani}, J., {Fabrizio}, M., {et~al.} 2021, \apj, 919, 85

\bibitem[{{Cioni} {et~al.}(2011){Cioni}, {Clementini}, {Girardi}, {Guandalini},
  {Gullieuszik}, {Miszalski}, {Moretti}, {Ripepi}, {Rubele}, {Bagheri},
  {Bekki}, {Cross}, {de Blok}, {de Grijs}, {Emerson}, {Evans}, {Gibson},
  {Gonzales-Solares}, {Groenewegen}, {Irwin}, {Ivanov}, {Lewis}, {Marconi},
  {Marquette}, {Mastropietro}, {Moore}, {Napiwotzki}, {Naylor}, {Oliveira},
  {Read}, {Sutorius}, {van Loon}, {Wilkinson}, \& {Wood}}]{cioni11}
{Cioni}, M.-R.~L., {Clementini}, G., {Girardi}, L., {et~al.} 2011, \aap, 527,
  A116

\bibitem[{{Coppola} {et~al.}(2015){Coppola}, {Marconi}, {Stetson}, {Bono},
  {Braga}, {Ripepi}, {Dall'Ora}, {Musella}, {Buonanno}, {Fabrizio}, {Ferraro},
  {Fiorentino}, {Iannicola}, {Monelli}, {Nonino}, {Th{\'e}venin}, \&
  {Walker}}]{coppola15}
{Coppola}, G., {Marconi}, M., {Stetson}, P.~B., {et~al.} 2015, \apj, 814, 71

\bibitem[{{Crestani} {et~al.}(2021){Crestani}, {Braga}, {Fabrizio}, {Bono},
  {Sneden}, {Preston}, {Ferraro}, {Iannicola}, {Nonino}, {Fiorentino},
  {Th{\'e}venin}, {Lemasle}, {Prudil}, {Alves-Brito}, {Altavilla}, {Chaboyer},
  {Dall'Ora}, {D'Orazi}, {Gilligan}, {Grebel}, {Koch-Hansen}, {Lala},
  {Marengo}, {Marinoni}, {Marrese}, {Mart{\'\i}nez-V{\'a}zquez}, {Matsunaga},
  {Monelli}, {Mullen}, {Neeley}, {da Silva}, {Stetson}, {Salaris}, {Storm},
  {Valenti}, \& {Zoccali}}]{crestani21b}
{Crestani}, J., {Braga}, V.~F., {Fabrizio}, M., {et~al.} 2021, \apj, 914, 10

\bibitem[{{Dall'Ora} {et~al.}(2003){Dall'Ora}, {Ripepi}, {Caputo},
  {Castellani}, {Bono}, {Smith}, {Brocato}, {Buonanno}, {Castellani}, {Corsi},
  {Marconi}, {Monelli}, {Nonino}, {Pulone}, \& {Walker}}]{dallora03}
{Dall'Ora}, M., {Ripepi}, V., {Caputo}, F., {et~al.} 2003, \aj, 126, 197

\bibitem[{{De Somma} {et~al.}(2020){De Somma}, {Marconi}, {Cassisi}, {Ripepi},
  {Leccia}, {Molinaro}, \& {Musella}}]{desomma21a}
{De Somma}, G., {Marconi}, M., {Cassisi}, S., {et~al.} 2020, \mnras, 496, 5039

\bibitem[{{Dekel} \& {Silk}(1986)}]{dekel86}
{Dekel}, A. \& {Silk}, J. 1986, \apj, 303, 39

\bibitem[{{Elorrieta} {et~al.}(2016){Elorrieta}, {Eyheramendy}, {Jord{\'a}n},
  {D{\'e}k{\'a}ny}, {Catelan}, {Angeloni}, {Alonso-Garc{\'\i}a},
  {Contreras-Ramos}, {Gran}, {Hajdu}, {Espinoza}, {Saito}, \&
  {Minniti}}]{elorrieta16}
{Elorrieta}, F., {Eyheramendy}, S., {Jord{\'a}n}, A., {et~al.} 2016, \aap, 595,
  A82

\bibitem[{{Feast} \& {Walker}(1987)}]{feast87}
{Feast}, M.~W. \& {Walker}, A.~R. 1987, \araa, 25, 345

\bibitem[{{Gallart} {et~al.}(2019){Gallart}, {Bernard}, {Brook}, {Ruiz-Lara},
  {Cassisi}, {Hill}, \& {Monelli}}]{gallart19}
{Gallart}, C., {Bernard}, E.~J., {Brook}, C.~B., {et~al.} 2019, Nature
  Astronomy, 3, 932

\bibitem[{{Gallart} {et~al.}(2005){Gallart}, {Zoccali}, \&
  {Aparicio}}]{gallart05}
{Gallart}, C., {Zoccali}, M., \& {Aparicio}, A. 2005, \araa, 43, 387

\bibitem[{{Hernitschek} {et~al.}(2016){Hernitschek}, {Schlafly}, {Sesar},
  {Rix}, {Hogg}, {Ivezi{\'c}}, {Grebel}, {Bell}, {Martin}, {Burgett},
  {Flewelling}, {Hodapp}, {Kaiser}, {Magnier}, {Metcalfe}, {Wainscoat}, \&
  {Waters}}]{hernitschek16}
{Hernitschek}, N., {Schlafly}, E.~F., {Sesar}, B., {et~al.} 2016, \apj, 817, 73

\bibitem[{{Hodge} \& {Wright}(1978)}]{hodge78}
{Hodge}, P.~W. \& {Wright}, F.~W. 1978, \aj, 83, 228

\bibitem[{{Holl} {et~al.}(2018){Holl}, {Audard}, {Nienartowicz}, {Jevardat de
  Fombelle}, {Marchal}, {Mowlavi}, {Clementini}, {De Ridder}, {Evans}, {Guy},
  {Lanzafame}, {Lebzelter}, {Rimoldini}, {Roelens}, {Zucker}, {Distefano},
  {Garofalo}, {Lecoeur-Ta{\"\i}bi}, {Lopez}, {Molinaro}, {Muraveva}, {Panahi},
  {Regibo}, {Ripepi}, {Sarro}, {Aerts}, {Anderson}, {Charnas}, {Barblan},
  {Blanco-Cuaresma}, {Busso}, {Cuypers}, {De Angeli}, {Glass}, {Grenon},
  {Juh{\'a}sz}, {Kochoska}, {Koubsky}, {Lanza}, {Leccia}, {Lorenz}, {Marconi},
  {Marschalk{\'o}}, {Mazeh}, {Messina}, {Mignard}, {Moitinho}, {Moln{\'a}r},
  {Morgenthaler}, {Musella}, {Ordenovic}, {Ord{\'o}{\~n}ez}, {Pagano},
  {Palaversa}, {Pawlak}, {Plachy}, {Pr{\v{s}}a}, {Riello}, {S{\"u}veges},
  {Szabados}, {Szegedi-Elek}, {Votruba}, \& {Eyer}}]{holl18}
{Holl}, B., {Audard}, M., {Nienartowicz}, K., {et~al.} 2018, \aap, 618, A30

\bibitem[{{Hosenie} {et~al.}(2019){Hosenie}, {Lyon}, {Stappers}, \&
  {Mootoovaloo}}]{hosenie19}
{Hosenie}, Z., {Lyon}, R.~J., {Stappers}, B.~W., \& {Mootoovaloo}, A. 2019,
  \mnras, 488, 4858

\bibitem[{{Kinemuchi} {et~al.}(2008){Kinemuchi}, {Harris}, {Smith},
  {Silbermann}, {Snyder}, {La Cluyz{\'e}}, \& {Clark}}]{kinemuchi08}
{Kinemuchi}, K., {Harris}, H.~C., {Smith}, H.~A., {et~al.} 2008, \aj, 136, 1921

\bibitem[{{Kolenberg} {et~al.}(2010){Kolenberg}, {Szab{\'o}}, {Kurtz},
  {Gilliland}, {Christensen-Dalsgaard}, {Kjeldsen}, {Brown}, {Benk{\H{o}}},
  {Chadid}, {Derekas}, {Di Criscienzo}, {Guggenberger}, {Kinemuchi}, {Kunder},
  {Koll{\'a}th}, {Kopacki}, {Moskalik}, {Nemec}, {Nuspl}, {Silvotti}, {Suran},
  {Borucki}, {Koch}, \& {Jenkins}}]{kolenberg10}
{Kolenberg}, K., {Szab{\'o}}, R., {Kurtz}, D.~W., {et~al.} 2010, \apjl, 713,
  L198

\bibitem[{{Marconi} {et~al.}(2015){Marconi}, {Coppola}, {Bono}, {Braga},
  {Pietrinferni}, {Buonanno}, {Castellani}, {Musella}, {Ripepi}, \&
  {Stellingwerf}}]{marconi15}
{Marconi}, M., {Coppola}, G., {Bono}, G., {et~al.} 2015, \apj, 808, 50

\bibitem[{{Mart{\'{\i}}nez-V{\'a}zquez}
  {et~al.}(2016){Mart{\'{\i}}nez-V{\'a}zquez}, {Monelli}, {Gallart}, {Bono},
  {Bernard}, {Stetson}, {Ferraro}, {Walker}, {Dall'Ora}, {Fiorentino}, \&
  {Iannicola}}]{martinezvazquez16a}
{Mart{\'{\i}}nez-V{\'a}zquez}, C.~E., {Monelli}, M., {Gallart}, C., {et~al.}
  2016, \mnras, 461, L41

\bibitem[{{Massari} {et~al.}(2019){Massari}, {Koppelman}, \&
  {Helmi}}]{massari19}
{Massari}, D., {Koppelman}, H.~H., \& {Helmi}, A. 2019, \aap, 630, L4

\bibitem[{{Mayer} {et~al.}(2001){Mayer}, {Governato}, {Colpi}, {Moore},
  {Quinn}, {Wadsley}, {Stadel}, \& {Lake}}]{mayer01c}
{Mayer}, L., {Governato}, F., {Colpi}, M., {et~al.} 2001, \apj, 559, 754

\bibitem[{{Nemec} {et~al.}(1988){Nemec}, {Wehlau}, \& {Mendes de
  Oliveira}}]{nemec88}
{Nemec}, J.~M., {Wehlau}, A., \& {Mendes de Oliveira}, C. 1988, \aj, 96, 528

\bibitem[{{Richards} {et~al.}(2011){Richards}, {Starr}, {Butler}, {Bloom},
  {Brewer}, {Crellin-Quick}, {Higgins}, {Kennedy}, \& {Rischard}}]{richards11}
{Richards}, J.~W., {Starr}, D.~L., {Butler}, N.~R., {et~al.} 2011, \apj, 733,
  10

\bibitem[{{Riess} {et~al.}(2021){Riess}, {Casertano}, {Yuan}, {Bowers},
  {Macri}, {Zinn}, \& {Scolnic}}]{riess21}
{Riess}, A.~G., {Casertano}, S., {Yuan}, W., {et~al.} 2021, \apjl, 908, L6

\bibitem[{{Sawala} {et~al.}(2010){Sawala}, {Scannapieco}, {Maio}, \&
  {White}}]{sawala10}
{Sawala}, T., {Scannapieco}, C., {Maio}, U., \& {White}, S. 2010, \mnras, 402,
  1599

\bibitem[{{Siegel} \& {Majewski}(2000)}]{siegel00}
{Siegel}, M.~H. \& {Majewski}, S.~R. 2000, \aj, 120, 284

\bibitem[{{Soszy{\'n}ski} {et~al.}(2008{\natexlab{a}}){Soszy{\'n}ski},
  {Poleski}, {Udalski}, {Szymanski}, {Kubiak}, {Pietrzynski}, {Wyrzykowski},
  {Szewczyk}, \& {Ulaczyk}}]{soszynski08b}
{Soszy{\'n}ski}, I., {Poleski}, R., {Udalski}, A., {et~al.} 2008{\natexlab{a}},
  \actaa, 58, 163

\bibitem[{{Soszy{\'n}ski} {et~al.}(2019){Soszy{\'n}ski}, {Smolec}, {Udalski},
  \& {Pietrukowicz}}]{soszynski19a}
{Soszy{\'n}ski}, I., {Smolec}, R., {Udalski}, A., \& {Pietrukowicz}, P. 2019,
  \apj, 873, 43

\bibitem[{{Soszy{\'n}ski} {et~al.}(2008{\natexlab{b}}){Soszy{\'n}ski},
  {Udalski}, {Szyma{\'n}ski}, {Kubiak}, {Pietrzy{\'n}ski}, {Wyrzykowski},
  {Szewczyk}, {Ulaczyk}, \& {Poleski}}]{soszynski08c}
{Soszy{\'n}ski}, I., {Udalski}, A., {Szyma{\'n}ski}, M.~K., {et~al.}
  2008{\natexlab{b}}, \actaa, 58, 293

\bibitem[{{Soszy{\'n}ski} {et~al.}(2009){Soszy{\'n}ski}, {Udalski},
  {Szyma{\'n}ski}, {Kubiak}, {Pietrzy{\'n}ski}, {Wyrzykowski}, {Szewczyk},
  {Ulaczyk}, \& {Poleski}}]{soszynski09a}
{Soszy{\'n}ski}, I., {Udalski}, A., {Szyma{\'n}ski}, M.~K., {et~al.} 2009,
  \actaa, 59, 1

\bibitem[{{Swope}(1967)}]{swope67}
{Swope}, H.~H. 1967, \pasp, 79, 439

\bibitem[{{Thackeray}(1950)}]{thackeray50}
{Thackeray}, A.~D. 1950, The Observatory, 70, 144

\bibitem[{{Thackeray}(1951)}]{thackeray51}
{Thackeray}, A.~D. 1951, The Observatory, 71, 219

\bibitem[{{Thackeray}(1958)}]{thackeray58}
{Thackeray}, A.~D. 1958, \mnras, 118, 117

\bibitem[{{Thackeray} \& {Wesselink}(1953)}]{thackeray53}
{Thackeray}, A.~D. \& {Wesselink}, A.~J. 1953, \nat, 171, 693

\bibitem[{{Tolstoy} {et~al.}(2009){Tolstoy}, {Hill}, \& {Tosi}}]{tolstoy09}
{Tolstoy}, E., {Hill}, V., \& {Tosi}, M. 2009, \araa, 47, 371

\bibitem[{{Udalski} {et~al.}(1992){Udalski}, {Szymanski}, {Kaluzny}, {Kubiak},
  \& {Mateo}}]{udalski92}
{Udalski}, A., {Szymanski}, M., {Kaluzny}, J., {Kubiak}, M., \& {Mateo}, M.
  1992, \actaa, 42, 253

\bibitem[{{van Agt}(1967)}]{vanagt67}
{van Agt}, S.~L. T.~J. 1967, \bain, 19, 275

\bibitem[{{van Agt}(1968)}]{vanagt68}
{van Agt}, S.~L. T.~J. 1968, Bulletin of the Astronomical Institutes of the
  Netherlands Supplement Series, 2, 237

\bibitem[{{Vivas} {et~al.}(2019){Vivas}, {Alonso-Garc{\'\i}a}, {Mateo},
  {Walker}, \& {Howard}}]{vivas19}
{Vivas}, A.~K., {Alonso-Garc{\'\i}a}, J., {Mateo}, M., {Walker}, A., \&
  {Howard}, B. 2019, \aj, 157, 35

\bibitem[{{Vivas} {et~al.}(2020){Vivas}, {Mart{\'\i}nez-V{\'a}zquez}, \&
  {Walker}}]{vivas20b}
{Vivas}, A.~K., {Mart{\'\i}nez-V{\'a}zquez}, C., \& {Walker}, A.~R. 2020,
  \apjs, 247, 35

\bibitem[{{Walker}(1989)}]{walker89}
{Walker}, A.~R. 1989, \aj, 98, 2086


\end{thebibliography}


\end{document}

